\section{Conclusions}
\label{rz:conclusiones}

On each fixed window, the Weil form admits an exact
reduction to a finite problem that preserves its
infinite-dimensional complement. The space of zero-mean
details can be made coercive by refinement; its elimination
produces a Schur complement with an equivalent positivity
criterion and the same negative index as the complete form. On the window
$(-1/18,1/18)$, joint control of averages, details and
couplings proves
\[
\mathscr W(f,f)\geq\frac12\|f\|_2^2,
\qquad f\in V_{1/9}.
\]
The assertion comprises the closure of smooth compactly
supported test functions in the joint readout norm.
It is a result on an infinite-dimensional functional
domain, stable under translation and interior refinement.

\subsection*{Sign of the form and preservation of the infinite complement}

The identity
\[
\mathscr W_R(Ua+d,Ua+d)
=a^*\bigl(A-B^*D^{-1}B\bigr)a
+\mathfrak d[d+D^{-1}Ba,d+D^{-1}Ba]
\]
explains what the reduction preserves. The second term
is positive and contains all details; the first retains
their interaction with the averages through $B^*D^{-1}B$.
Theorem~\ref{schur-add-schur-exacto} proves that positivity
on $V_R$ is equivalent to positivity of this matrix.
Its finite size does not replace calculation of its
entries: the residual certification in
\eqref{schur-add-eq-S7} supplies bounds retaining the
complement's error.

On the window of length $1/9$, the bounds
$D>I$, $A>97/100$ and $\|B\|^2<1/5$ produce the stated
coercive margin. This local case contains no active prime
translation because $1/9<\log2$. Invariance of the
complete form nevertheless carries the same bound to
each subspace $C_a^n\tau_t C_c^\infty((-1/18,1/18))$,
with $a>b>1/2$. Its elements may have exponential tails,
and their prime contribution is preserved with the full
convergent sum. The estimate applies to each entire
collection, not merely to a selected test function.

The other positive domains retain their own content:
the specified assembly of nine intervals controls its
cross products and all refinements, while an arbitrary
window admits a coercivity threshold for cell details.
Distinguishing these domains identifies both the proved
margin and the coupling determining the sign of a
broader composition.

\subsection*{Arithmetic determines the modes and their readouts}

The arithmetic antecedent is constructed before spectral
readout. The two APP sheets supply the local operations
and their quotients; the native product propagates these
data through words with unique normal forms. Product
decompositions select irreducibles, and evaluation
identifies them with positive primes. The operator
realization of the selector counts nontrivial
decompositions. Its positivity is a squared norm with
this concrete arithmetic meaning.

Factorization subsequently organizes mode occupations
and their clocks. The Euler product and prime--power
moments arise from traces of that representation.
The limit of the cyclic heat trace and its Mellin
transform produce the completed function, with its
gamma and polar factors and functional equation.
Finite parts, regularized evaluations and remainder
bounds preserve the normalizations required to compose
these readouts.

The unity of this route lies in the enriched
APP--TRIT--TPK states and the joint structure of the
continuum: route, orientation, sheets, memory and
incidences are transmitted through compatible
restrictions and refinements. Real evaluation obtains
a result of this construction; history and boundary
retain the information needed by its subsequent
realizations.

The composition laws make this continuity precise. The product of two
moments is realized on pairs of modes, and its coefficients are grouped
by divisor convolution; diagonal composition retains a single mode.
Spectral normalization transports the entire Laurent expansion,
including the polar term, through an invertible triangular affine
transformation. The functional equation relates negative even
derivatives to positive odd values and, in particular, establishes
$Z(3)=4\pi^2\log A_2$. Finally, irreducible-mode cutoffs provide
explicit enclosures complementary to depth cutoffs. Each identity
retains the operation producing it and the domain in which it is proved.

\subsection*{Boundary, vacancies and double circle}

Reconstruction of the interior from the lifted boundary
in the multiaffine chart gives an explicit instance of
preservation. The sum--product difference retains both
the residual datum and the quotient; the readouts
$72,27,77,22$ and the completion $729+271$ retain their
operations and domains. The center $(5,5)$ and its
conjugate vacancies preserve sum and multiplicative
residue while the quotient changes. The regional
signature $555555$ requires three continuation sheets,
and the capacity calendar produces lengths $21,22$
and separations $6,7$ with an exact balance.

The double circle carries this preservation into
readout transport. Nine-phase compression produces
the complement $(8z,-z,\ldots,-z)/27$ and the factor
$8/81$. The modal ratio $q=5/9$ determines the memory
moments $q^{|n|}$; their triangular factorization proves
positivity in every dimension and the determinant
$(56/81)^N$. The loss register preserves the terminal
and cross products. Refinement may lead to a weak
atomic limit of states, whose representation is
distinguished from the regular shift in which the
initial vectors were constructed.

For the arithmetic readouts, the gamma tail
reconstructs the test function and the rows of both
columns. The resulting connector preserves exactly
$J_+^*J_+-J_-^*J_-$, with the prime, gamma and polar
terms. Its action formulates the energy problem on
an explicit operator. Invariance under $C_a$ also
preserves all cross moments of the form, not only
its diagonal evaluations.

\subsection*{Global condition and relation to the Riemann hypothesis}

Weil's global criterion requires positivity on the
entire test-function domain. In the two-column
formulation, the exact condition is contractivity of
the reachable transfer on the closure of the positive
image, with compatibility as windows enlarge. In the
circular formulation, it is expressed by the identity
\[
m_n^{\rm ar}(f,g)
=\langle U_{\rm mem}^{\,n}\mathcal E f,
          \mathcal E g\rangle_{\mathcal K_{\rm lim}},
\]
for all mixed moments and a collection covering the
domain through the stated limits. The energy identity
and intertwining must hold jointly; positivity of
the regular memory Gram form does not replace them.

The proofs in this article establish reconstruction,
invariance and coercivity on the stated domains.
The global identification above remains a hypothesis
of the spectral results employing it and does not
follow from positivity of each transported collection
separately. The article therefore does not conclude
a resolution of the Riemann hypothesis. Its functional
contribution is precise: preservation of arithmetic
structure and memory is expressed through recoverable
operators, exact reductions and coercive bounds,
with the mixed products and domains governing any
global extension kept explicit.
